\documentclass[11pt,a4paper]{article}
\usepackage[utf8]{inputenc}
\usepackage[T1]{fontenc}
\usepackage{lmodern}
\usepackage[french]{babel}
\usepackage{amsmath,amssymb,mathtools}
\usepackage{icomma}
\usepackage{xcolor}
\usepackage{colortbl}
\usepackage{tikz}
\usepackage[most]{tcolorbox}
\usepackage{enumitem}
\usepackage{array}
\usepackage{booktabs}
\usepackage{fancyhdr}
\usepackage[hidelinks]{hyperref}
\usetikzlibrary{arrows.meta,positioning,calc,patterns,backgrounds,shapes.geometric,decorations.pathreplacing,decorations.markings}
\usepackage[a4paper,margin=2.15cm,headheight=26pt,headsep=12pt]{geometry}
\usepackage{titlesec}

% ---------------- Palette ----------------
\definecolor{primary}{RGB}{0,151,169}
\definecolor{primarydark}{RGB}{0,88,101}
\definecolor{defcol}{RGB}{0,114,178}
\definecolor{thmcol}{RGB}{0,140,120}
\definecolor{formulacol}{RGB}{198,93,9}
\definecolor{excol}{RGB}{0,135,90}
\definecolor{attncol}{RGB}{200,40,55}
\definecolor{methcol}{RGB}{90,90,100}
\definecolor{remarkcol}{RGB}{120,94,200}
\definecolor{barcol}{RGB}{0,151,169}
\definecolor{barcold}{RGB}{0,110,124}
\definecolor{modecol}{RGB}{200,55,55}
\definecolor{meancol}{RGB}{0,90,200}
\definecolor{medcol}{RGB}{160,90,190}
\definecolor{quartcol}{RGB}{0,135,90}
\definecolor{ogivecol}{RGB}{176,74,0}

% ---------------- Boîtes ----------------
\tcbset{boxbase/.style={enhanced,breakable,boxrule=0.9pt,arc=2.5pt,
  left=3mm,right=3mm,top=2.3mm,bottom=2.3mm,
  fonttitle=\bfseries,coltitle=white,toptitle=0.6mm,bottomtitle=0.6mm}}

\newtcolorbox{definition}[1][]{boxbase,colback=defcol!6,colframe=defcol,
  title={\textsc{Définition}\ifx\relax#1\relax\relax\else\textmd{ : \itshape #1}\fi}}
\newtcolorbox{propriete}[1][]{boxbase,colback=thmcol!7,colframe=thmcol,
  title={\textsc{Propriété}\ifx\relax#1\relax\relax\else\textmd{ : \itshape #1}\fi}}
\newtcolorbox{formule}[1][]{boxbase,colback=formulacol!7,colframe=formulacol,
  title={\textsc{Formule}\ifx\relax#1\relax\relax\else\textmd{ : \itshape #1}\fi}}
\newtcolorbox{exemple}[1][]{boxbase,colback=excol!6,colframe=excol,
  title={\textsc{Exemple}\ifx\relax#1\relax\relax\else\textmd{ : \itshape #1}\fi}}
\newtcolorbox{methode}[1][]{boxbase,colback=methcol!8,colframe=methcol,sharp corners,
  title={\textsc{Méthode}\ifx\relax#1\relax\relax\else\textmd{ : \itshape #1}\fi}}
\newtcolorbox{remarque}[1][]{boxbase,colback=remarkcol!7,colframe=remarkcol,
  title={\textsc{Astuce}\ifx\relax#1\relax\relax\else\textmd{ : \itshape #1}\fi}}
\newtcolorbox{attention}[1][]{boxbase,colback=attncol!7,colframe=attncol,
  title={\textsc{Attention}\ifx\relax#1\relax\relax\else\textmd{ : \itshape #1}\fi}}

% Boîte "exercice" et "corrigé" (titre libre passé en argument)
\newtcolorbox{exobox}[1]{boxbase,colback=primary!4,colframe=primary,
  before skip=8pt,after skip=8pt,title={#1}}
\newtcolorbox{corbox}[1]{boxbase,colback=excol!5,colframe=excol!85!black,
  before skip=8pt,after skip=8pt,title={#1}}

% ---------------- En-tête ----------------
\newcommand{\docpart}[1]{\def\thedocpart{#1}}
\docpart{}
\pagestyle{fancy}
\fancyhf{}
\fancyhead[L]{\small\textcolor{primarydark}{\textbf{Fiche 8} — Statistiques descriptives}}
\fancyhead[R]{\small\textcolor{primarydark}{\textsc{\thedocpart}}}
\fancyfoot[C]{\small\thepage}
\renewcommand{\headrulewidth}{0.8pt}
\renewcommand{\headrule}{\hbox to\headwidth{\color{primary}\leaders\hrule height \headrulewidth\hfill}}

\titleformat{\section}{\large\bfseries\color{primarydark}}{\thesection}{0.6em}{}
\titleformat{\subsection}{\normalsize\bfseries\color{primary}}{\thesubsection}{0.6em}{}
\setlist{topsep=2pt,itemsep=1.5pt,parsep=0pt}

% Bandeau de titre du document
\newcommand{\bandeau}[2]{%
\begin{tcolorbox}[boxbase,colback=primary,colframe=primarydark,halign=center,
   before skip=2pt,after skip=12pt]
{\LARGE\bfseries\color{white} #1}\\[3pt]{\normalsize\color{white} #2}
\end{tcolorbox}}

\newcommand{\ecm}{\sigma_m}
\newcommand{\ect}{\sigma}
\newcommand{\med}{M\!e}
\docpart{Tutoriel — Thème 3}
\begin{document}
\bandeau{Thème 3 — Médiane et quartiles}{Le partage de l'effectif et l'interpolation linéaire}

\noindent\textbf{But du thème.} La \textbf{médiane} coupe la série en deux moitiés égales ; les
\textbf{quartiles} la coupent en quatre. Deux situations, deux techniques : la série \emph{discrète}
(on trie et on compte) et la série \emph{en classes} (on interpole). L'interpolation linéaire est
\emph{la} mécanique-signature de cette fiche : on la répète pour la médiane, $Q_1$ et $Q_3$.

\section{Médiane d'une série discrète : trier puis compter}

\begin{formule}[médiane de $n$ valeurs rangées par ordre croissant]
\begin{itemize}
  \item $n$ \textbf{impair} : la médiane est la valeur de rang $\dfrac{n+1}{2}$ ;
  \item $n$ \textbf{pair} : la médiane est la moyenne des valeurs de rangs $\dfrac{n}{2}$ et
        $\dfrac{n}{2}+1$.
\end{itemize}
\end{formule}

\begin{attention}[réflexe nº1 : trier]
Tant que la série n'est pas \textbf{rangée par ordre croissant}, la médiane n'a aucun sens.
Le tri est toujours la première étape.
\end{attention}

\begin{exemple}[$n$ impair puis $n$ pair]
$\{7,2,9,4,5\}$ trié donne $2,4,\underline{5},7,9$ : $n=5$ impair, rang $\tfrac{5+1}{2}=3$, d'où
$\med=\mathbf{5}$.\\
$\{2,4,5,7,9,10\}$ : $n=6$ pair, rangs $3$ et $4$, d'où $\med=\dfrac{5+7}{2}=\mathbf{6}$.
\end{exemple}

\section{Médiane d'une série en classes : l'interpolation linéaire}

Avec des classes, on ne connaît plus les valeurs individuelles ; on suppose les individus
\textbf{répartis uniformément} dans la classe qui contient le milieu de l'effectif.

\begin{formule}[médiane par interpolation (triangles semblables)]
\[
\boxed{\;\med=a+\frac{r}{n_{\text{classe}}}\,(b-a)\;}\qquad r=\frac{N}{2}-n_c^{\text{préc}}
\]
\begin{itemize}
  \item $[a,b[$ : la \textbf{classe médiane} = première classe dont l'effectif cumulé atteint/dépasse
        $\tfrac{N}{2}$ ;
  \item $n_{\text{classe}}$ : effectif de cette classe ; $n_c^{\text{préc}}$ : cumul de la classe
        \emph{précédente} ;
  \item $r$ : rang du $\left(\tfrac{N}{2}\right)^{\text{e}}$ individu \emph{à l'intérieur} de la classe.
\end{itemize}
\end{formule}

\begin{methode}[4 étapes, valables aussi pour $Q_1$ et $Q_3$]
\begin{enumerate}
  \item \textbf{Rang visé} : $\tfrac{N}{2}$ (médiane), $\tfrac{N}{4}$ ($Q_1$), $\tfrac{3N}{4}$ ($Q_3$).
  \item \textbf{Localiser la classe} avec les effectifs cumulés $n_c$.
  \item \textbf{Rang intérieur} $r=(\text{rang visé})-n_c^{\text{préc}}$.
  \item \textbf{Interpoler} : (valeur)$=a+\dfrac{r}{n_{\text{classe}}}(b-a)$.
\end{enumerate}
\end{methode}

\begin{exemple}[médiane par interpolation]
Série ($N=66$) : cumuls $\ldots,30$ (classe $[85,95[$), puis $44$ (classe $[95,105[$).
Rang visé $\tfrac{66}{2}=33$. La classe médiane est $[95,105[$ ($a=95,b=105,n_{\text{classe}}=14$),
et $r=33-30=3$ :
\[ \med=95+\frac{3}{14}\times(105-95)=95+\frac{30}{14}=95+2{,}14=\mathbf{97{,}14}. \]
Interprétation : $50\,\%$ des individus sont sous $97{,}14$.
\end{exemple}

\section{Les quartiles}

\begin{definition}[quartiles]
$Q_1$ (rang $\tfrac{N}{4}$) : $25\,\%$ des données lui sont inférieures. $Q_3$ (rang $\tfrac{3N}{4}$) :
$75\,\%$. Et $Q_2=\med$. L'\textbf{écart interquartile} $Q_3-Q_1$ mesure la dispersion des $50\,\%$
« centraux », peu sensible aux extrêmes.
\end{definition}

\begin{remarque}[même moule pour tout]
Médiane, $Q_1$, $Q_3$ : \emph{exactement} la même formule d'interpolation, seul le \textbf{rang visé}
change ($\tfrac N2$, $\tfrac N4$, $\tfrac{3N}4$). Maîtriser un cas, c'est les maîtriser tous.
\end{remarque}
\end{document}
